%!TeX root=../main.tex

\chapter{Introdução}
\label{cap:introducao}

Em 20 de Abril de 1951 começa oficialmente a operação do \textit{Whirlwind I}, o primeiro computador
a conseguir exibir imagens gráficas em um monitor de tubo de raios catódicos em tempo real,
permitindo interação com o usuário. Desenvolvido por Jay Forrester e Robert Everett no
MIT (\textit{Massachusetts Institute of Technology}) e financiado pelas Forças Armadas dos Estados
Unidos, este computador tinha como objetivo inicial ser um simulador de voo analógico, mas
que eventualmente evoluiu para um sistema eletrônico digital. Embora desenvolvido para este
propósito, este computador já contava com um dos primeiros jogos eletrônicos da história, o
\textit{bouncing ball game}, feito para demonstrar a simulação física de movimento de uma bola quicando.

Desde então, a computação gráfica e os jogos evoluíram lado a lado, impulsionados por uma
busca contínua pelo fotorrealismo e por uma simulação fiel do mundo real. Essa corrida
tecnológica existe até os dias de hoje, com o exemplo recente mais famoso sendo a
implementação do DLSS5 em placas de vídeo modernas, uma tecnologia que utiliza redes neurais
e algoritmos de aceleração baseados em Inteligência Artificial para renderizar iluminação em
tempo real e cenas geometricamente complexas.

No entanto, para atingir este nível de processamento, a computação gráfica foi limitada pelas
restrições de hardware ao decorrer da história. Nas décadas de 1980 e 1990, a incapacidade de
processar grandes volumes de dados em tempo real exigia soluções criativas de otimização,
sendo então consolidadas técnicas como a modelagem de baixa poligonização (\textit{low-poly}), para
uma renderização mais eficiente, e a projeção isométrica, utilizada para simular profundidade e
exibir grandes cenários tridimensionais em ambientes 2D com baixo custo de processamento.

Com o avanço tecnológico, técnicas como essas deixaram de ser requisitos técnicos e se
tornaram uma linguagem de expressão visual. Enquanto jogos de alto orçamento (\textit{AAA})
continuam perseguindo a complexidade gráfica, desenvolvedores independentes conseguem
usar desta linguagem para liberar recursos de hardware e focar em outros aspectos de seus
jogos, como um desempenho mais acessível para o público e em como os aspectos visuais
podem afetar o \textit{game design} e o \textit{level design} do projeto.

Nesse quesito, o estilo do jogo influencia diretamente o \textit{game design} --- conjunto de regras,
mecânicas e dinâmicas de interação --- e o \textit{level design}, responsável pelo planejamento
espacial, arquitetura dos ambientes e pela curva de desafio oferecida ao jogador. Para este
projeto, destaca-se em particular jogos das categorias de aventura e \textit{fast-paced}. Enquanto
jogos de aventura exigem cenários visualmente desenvolvidos que estimulem a exploração, a
navegação e a resolução de quebra-cabeças, jogos \textit{fast-paced} demandam uma movimentação
dinâmica e precisa, com respostas imediatas aos comandos do jogador.

Portanto, com a redução de detalhes em modelos \textit{low-poly} e a projeção isométrica da câmera
do jogo, é possível manipular como o jogador identifica obstáculos, trajetórias e diferentes
geometrias de cenário enquanto executa comandos, navega e enfrenta diferentes desafios.

Então, sob a perspectiva deste trabalho, a escolha destas técnicas permite a exploração das
etapas do \textit{pipeline} gráfico para aprofundar o conhecimento em \textbf{Aplicações de Computação Gráfica
na construção de visuais estilizados, \textit{level design} e jogabilidade para um jogo 3D}, por
meio da implementação de diferentes interações entre objetos e métodos de renderização de
modelos, pixels e iluminação.

\section**{Objetivos}

O objetivo deste trabalho é projetar, desenvolver e documentar um jogo eletrônico de aventura
e ação \textit{fast-paced} em perspectiva isométrica 3D, utilizando estética \textit{low-poly} pixelizada, com
foco na integração entre mecânicas de física, computação gráfica, \textit{game design} e \textit{level design}.

Para atingir este objetivo, foram estabelecidas as seguintes metas:

\begin{description}
	\item[Desenvolver a arquitetura básica de movimentação e colisão:] Implementar um sistema
	de movimentação dinâmico e responsivo (\textit{fast-paced}) focado em física, incluindo
	mecânicas de aceleração, deslize em superfícies e navegação fluida em rampas e
	obstáculos;
	\item[Estruturar o primeiro cenário (Movimentação):] Projetar desafios de \textit{level design}
	baseados em navegação avançada e mecânicas de mobilidade (\textit{parkour});
	\item[Estruturar o segundo cenário (Interação):] Implementar sistemas de manipulação e
	destruição de objetos tridimensionais do cenário, permitindo a abertura de novas rotas e
	resolução de obstáculos;
	\item[Estruturar o terceiro cenário (Iluminação):] Desenvolver quebra-cabeças e dinâmicas de
	ambiente fundamentados em técnicas de iluminação e reatividade de materiais à luz;
	\item[Implementar o confronto final e integração de mecânicas:] Projetar um combate contra
	um chefe final que exija a aplicação combinada de todas as ferramentas e mecânicas
	adquiridas ao longo dos cenários anteriores;
	\item[Aplicar técnicas de estilização no \textit{pipeline} gráfico:] Integrar os modelos \textit{low-poly}, a
	renderização pixelizada e a projeção ortográfica isométrica, mantendo a legibilidade
	espacial e o desempenho do jogo.

\end{description}

\section**{Estrutura do Texto}

